% --- Archon physics formalization source begin ---
% archon:physics
% archon:covers PhyXMiniProblems/problem_phyx_mini_0815.lean
% archon:source-report reports/phyx_mini/problem_phyx_mini_0815.source.json

\chapter{Physics problem phyx\_mini\_0815}
\label{ch:PhyXMiniProblems_problem_phyx_mini_0815}

\paragraph{Problem source.}
\#\# Physical scenario

Figure shows a ballistic pendulum, a simple system for measuring the speed of a bullet. A bullet of mass \textbackslash{}( m\_\{\textbackslash{}text\{B\}\} \textbackslash{}) makes a completely inelastic collision with a block of wood of mass \textbackslash{}( m\_\{\textbackslash{}text\{W\}\} \textbackslash{}), which is suspended like a pendulum. After the impact, the block swings up to a maximum height \textbackslash{}( h \textbackslash{}). In terms of \textbackslash{}( h \textbackslash{}), \textbackslash{}( m\_\{\textbackslash{}text\{B\}\} \textbackslash{}), and \textbackslash{}( m\_\{\textbackslash{}text\{W\}\} \textbackslash{}).

\#\# Question

What is the initial speed \textbackslash{}( v\_1 \textbackslash{}) of the bullet?

\#\# Answer choices

- A: A: \textbackslash{}( \textbackslash{}frac\{m\_\{\textbackslash{}text\{B\}\} + m\_\{\textbackslash{}text\{W\}\}\}\{m\_\{\textbackslash{}text\{B\}\}\} \textbackslash{}sqrt\{gh\} \textbackslash{})

- B: B: \textbackslash{}( \textbackslash{}frac\{m\_\{\textbackslash{}text\{B\}\} + m\_\{\textbackslash{}text\{W\}\}\}\{m\_\{\textbackslash{}text\{B\}\}\} \textbackslash{}sqrt\{2gh\} \textbackslash{})

- C: C: \textbackslash{}( \textbackslash{}frac\{m\_\{\textbackslash{}text\{B\}\} + m\_\{\textbackslash{}text\{W\}\}\}\{m\_\{\textbackslash{}text\{W\}\}\} \textbackslash{}sqrt\{2gh\} \textbackslash{})

- D: D: \textbackslash{}( \textbackslash{}frac\{m\_\{\textbackslash{}text\{B\}\} + m\_\{\textbackslash{}text\{W\}\}\}\{m\_\{\textbackslash{}text\{W\}\}\} \textbackslash{}sqrt\{gh\} \textbackslash{})

\#\# Figure caption (auxiliary; use the image as primary evidence)

The image is divided into two sections, illustrating a physics scenario involving a collision and subsequent motion.

\#\#\# Top Section: Before Collision
- **Objects:**
  - A bullet labeled as \textbackslash{}( m\_B \textbackslash{}) is moving horizontally.
  - A wooden block labeled as \textbackslash{}( m\_W \textbackslash{}) is suspended by wires.
- **Motion:**
  - The bullet is moving towards the block with velocity \textbackslash{}( v\_1 \textbackslash{}).
- **Setup:**
  - There is a coordinate system with axes labeled \textbackslash{}( x \textbackslash{}) and \textbackslash{}( y \textbackslash{}).
- **Text:**
  - "Before collision" describes this scenario.

\#\#\# Bottom Section: Immediately After Collision
- **Objects:**
  - The bullet is embedded in the wooden block, now combined into a single object labeled \textbackslash{}( m\_B + m\_W \textbackslash{}).
- **Motion:**
  - The combined object moves with velocity \textbackslash{}( v\_2 \textbackslash{}).
  - The object swings upwards, reaching a height \textbackslash{}( h \textbackslash{}).
- **Setup:**
  - The original and final positions of the block are shown, indicating the arc of the swing.
- **Text:**
  - "Immediately after collision" describes this scenario.
  - "Top of swing" with an arrow indicates the maximum height reached.
  
Together, the diagrams illustrate a bullet-block collision, followed by the motion of the combined system.

\#\# Dataset metadata

Momentum and Collisions; Multi-Formula Reasoning, Physical Model Grounding Reasoning

\paragraph{Recorded answer/context.}
B: B: \textbackslash{}( \textbackslash{}frac\{m\_\{\textbackslash{}text\{B\}\} + m\_\{\textbackslash{}text\{W\}\}\}\{m\_\{\textbackslash{}text\{B\}\}\} \textbackslash{}sqrt\{2gh\} \textbackslash{})

\paragraph{Figure/image path.}
/root/proposal\_for\_physic/science-mango/phyx\_mini\_run/phyx\_data/test\_image/815.png

\paragraph{Formalization target.}
create a compiling Lean file with sorry bodies at `PhyXMiniProblems/problem\_phyx\_mini\_0815.lean`. The Lean declarations must preserve the physical quantities, dimensions or dimensional roles, figure labels, governing-law hypotheses, and final relation expressed by this problem.
Use Mathlib/Physlib names found through LeanExplore where available. If a physics API is missing, introduce faithful local abstractions rather than scalar placeholder aliases.

\begin{theorem}[Physics formalization target]
\label{thm:physics:phyx_mini_0815:target}
\lean{PhyXMiniProblems.ProblemPhyXMini0815.problem_phyx_mini_0815}
\uses{def:physics:phyx-mini-0815:phyxminiproblems-problemphyxmini0815-massinkilograms, def:physics:phyx-mini-0815:phyxminiproblems-problemphyxmini0815-lengthinmeters, def:physics:phyx-mini-0815:phyxminiproblems-problemphyxmini0815-speedinmeterspersecond, def:physics:phyx-mini-0815:phyxminiproblems-problemphyxmini0815-accelerationinmeterspersecondsquared, def:physics:phyx-mini-0815:phyxminiproblems-problemphyxmini0815-ballisticpendulumsetup, def:physics:phyx-mini-0815:phyxminiproblems-problemphyxmini0815-matchesstatedscenario, def:physics:phyx-mini-0815:phyxminiproblems-problemphyxmini0815-matchesprimaryfigure, def:physics:phyx-mini-0815:phyxminiproblems-problemphyxmini0815-hasphysicalparameters, def:physics:phyx-mini-0815:phyxminiproblems-problemphyxmini0815-satisfiesstickingcollisionmomentumlaws, def:physics:phyx-mini-0815:phyxminiproblems-problemphyxmini0815-satisfiesconservativependulumswinglaws}
The assigned autoformalize agent should translate this physics problem into Lean declarations in the covered file, with theorem and lemma proof bodies written as `by sorry`.
\end{theorem}
\begin{proof}
This is an autoformalization task, not a proof task. Produce statements that can later be proved by the physics prover without weakening the physical model.
\end{proof}

% --- Archon declaration topology begin ---
\section{Declaration topology}
\begin{definition}[acceleration Dimension]
  \label{def:physics:phyx-mini-0815:phyxminiproblems-problemphyxmini0815-accelerationdimension}
  \lean{PhyXMiniProblems.ProblemPhyXMini0815.accelerationDimension}
  Gravitational acceleration has dimension length divided by time squared
\end{definition}
\begin{proof}
  This is the declared model definition of acceleration Dimension.
\end{proof}
\begin{definition}[Mass Quantity]
  \label{def:physics:phyx-mini-0815:phyxminiproblems-problemphyxmini0815-massquantity}
  \lean{PhyXMiniProblems.ProblemPhyXMini0815.MassQuantity}
  A nonnegative, unit-independent physical mass
\end{definition}
\begin{proof}
  This is the declared model definition of Mass Quantity.
\end{proof}
\begin{definition}[Length Quantity]
  \label{def:physics:phyx-mini-0815:phyxminiproblems-problemphyxmini0815-lengthquantity}
  \lean{PhyXMiniProblems.ProblemPhyXMini0815.LengthQuantity}
  A nonnegative, unit-independent vertical height
\end{definition}
\begin{proof}
  This is the declared model definition of Length Quantity.
\end{proof}
\begin{definition}[Speed Quantity]
  \label{def:physics:phyx-mini-0815:phyxminiproblems-problemphyxmini0815-speedquantity}
  \lean{PhyXMiniProblems.ProblemPhyXMini0815.SpeedQuantity}
  A nonnegative, unit-independent speed magnitude
\end{definition}
\begin{proof}
  This is the declared model definition of Speed Quantity.
\end{proof}
\begin{definition}[Acceleration Quantity]
  \label{def:physics:phyx-mini-0815:phyxminiproblems-problemphyxmini0815-accelerationquantity}
  \lean{PhyXMiniProblems.ProblemPhyXMini0815.AccelerationQuantity}
  \uses{def:physics:phyx-mini-0815:phyxminiproblems-problemphyxmini0815-accelerationdimension}
  A nonnegative, unit-independent acceleration magnitude
\end{definition}
\begin{proof}
  This is the declared model definition of Acceleration Quantity.
\end{proof}
\begin{definition}[Horizontal Momentum Quantity]
  \label{def:physics:phyx-mini-0815:phyxminiproblems-problemphyxmini0815-horizontalmomentumquantity}
  \lean{PhyXMiniProblems.ProblemPhyXMini0815.HorizontalMomentumQuantity}
  A signed one-dimensional momentum component along the displayed x-axis
\end{definition}
\begin{proof}
  This is the declared model definition of Horizontal Momentum Quantity.
\end{proof}
\begin{definition}[Energy Quantity]
  \label{def:physics:phyx-mini-0815:phyxminiproblems-problemphyxmini0815-energyquantity}
  \lean{PhyXMiniProblems.ProblemPhyXMini0815.EnergyQuantity}
  Mechanical energy with Physlib's energy dimension
\end{definition}
\begin{proof}
  This is the declared model definition of Energy Quantity.
\end{proof}
\begin{definition}[mass Readout]
  \label{def:physics:phyx-mini-0815:phyxminiproblems-problemphyxmini0815-massreadout}
  \lean{PhyXMiniProblems.ProblemPhyXMini0815.massReadout}
  \uses{def:physics:phyx-mini-0815:phyxminiproblems-problemphyxmini0815-massquantity}
  Read a physical mass in a selected mass unit
\end{definition}
\begin{proof}
  This is the declared model definition of mass Readout.
\end{proof}
\begin{definition}[length Readout]
  \label{def:physics:phyx-mini-0815:phyxminiproblems-problemphyxmini0815-lengthreadout}
  \lean{PhyXMiniProblems.ProblemPhyXMini0815.lengthReadout}
  \uses{def:physics:phyx-mini-0815:phyxminiproblems-problemphyxmini0815-lengthquantity}
  Read a physical length in a selected length unit
\end{definition}
\begin{proof}
  This is the declared model definition of length Readout.
\end{proof}
\begin{definition}[speed Readout]
  \label{def:physics:phyx-mini-0815:phyxminiproblems-problemphyxmini0815-speedreadout}
  \lean{PhyXMiniProblems.ProblemPhyXMini0815.speedReadout}
  \uses{def:physics:phyx-mini-0815:phyxminiproblems-problemphyxmini0815-speedquantity}
  Read a speed in coherent selected length and time units
\end{definition}
\begin{proof}
  This is the declared model definition of speed Readout.
\end{proof}
\begin{definition}[acceleration Readout]
  \label{def:physics:phyx-mini-0815:phyxminiproblems-problemphyxmini0815-accelerationreadout}
  \lean{PhyXMiniProblems.ProblemPhyXMini0815.accelerationReadout}
  \uses{def:physics:phyx-mini-0815:phyxminiproblems-problemphyxmini0815-accelerationquantity}
  Read acceleration in coherent selected length and time units
\end{definition}
\begin{proof}
  This is the declared model definition of acceleration Readout.
\end{proof}
\begin{definition}[horizontal Momentum Readout]
  \label{def:physics:phyx-mini-0815:phyxminiproblems-problemphyxmini0815-horizontalmomentumreadout}
  \lean{PhyXMiniProblems.ProblemPhyXMini0815.horizontalMomentumReadout}
  \uses{def:physics:phyx-mini-0815:phyxminiproblems-problemphyxmini0815-horizontalmomentumquantity}
  Read the unique component of a one-dimensional momentum quantity
\end{definition}
\begin{proof}
  This is the declared model definition of horizontal Momentum Readout.
\end{proof}
\begin{definition}[energy Readout]
  \label{def:physics:phyx-mini-0815:phyxminiproblems-problemphyxmini0815-energyreadout}
  \lean{PhyXMiniProblems.ProblemPhyXMini0815.energyReadout}
  \uses{def:physics:phyx-mini-0815:phyxminiproblems-problemphyxmini0815-energyquantity}
  Read energy in the coherent unit induced by selected mechanical units
\end{definition}
\begin{proof}
  This is the declared model definition of energy Readout.
\end{proof}
\begin{definition}[mass In Kilograms]
  \label{def:physics:phyx-mini-0815:phyxminiproblems-problemphyxmini0815-massinkilograms}
  \lean{PhyXMiniProblems.ProblemPhyXMini0815.massInKilograms}
  \uses{def:physics:phyx-mini-0815:phyxminiproblems-problemphyxmini0815-massquantity, def:physics:phyx-mini-0815:phyxminiproblems-problemphyxmini0815-massreadout}
  Kilogram readout of a physical mass
\end{definition}
\begin{proof}
  This is the declared model definition of mass In Kilograms.
\end{proof}
\begin{definition}[length In Meters]
  \label{def:physics:phyx-mini-0815:phyxminiproblems-problemphyxmini0815-lengthinmeters}
  \lean{PhyXMiniProblems.ProblemPhyXMini0815.lengthInMeters}
  \uses{def:physics:phyx-mini-0815:phyxminiproblems-problemphyxmini0815-lengthquantity, def:physics:phyx-mini-0815:phyxminiproblems-problemphyxmini0815-lengthreadout}
  Metre readout of a physical height
\end{definition}
\begin{proof}
  This is the declared model definition of length In Meters.
\end{proof}
\begin{definition}[speed In Meters Per Second]
  \label{def:physics:phyx-mini-0815:phyxminiproblems-problemphyxmini0815-speedinmeterspersecond}
  \lean{PhyXMiniProblems.ProblemPhyXMini0815.speedInMetersPerSecond}
  \uses{def:physics:phyx-mini-0815:phyxminiproblems-problemphyxmini0815-speedquantity, def:physics:phyx-mini-0815:phyxminiproblems-problemphyxmini0815-speedreadout}
  Metre-per-second readout of a speed magnitude
\end{definition}
\begin{proof}
  This is the declared model definition of speed In Meters Per Second.
\end{proof}
\begin{definition}[acceleration In Meters Per Second Squared]
  \label{def:physics:phyx-mini-0815:phyxminiproblems-problemphyxmini0815-accelerationinmeterspersecondsquared}
  \lean{PhyXMiniProblems.ProblemPhyXMini0815.accelerationInMetersPerSecondSquared}
  \uses{def:physics:phyx-mini-0815:phyxminiproblems-problemphyxmini0815-accelerationquantity, def:physics:phyx-mini-0815:phyxminiproblems-problemphyxmini0815-accelerationreadout}
  Metre-per-second-squared readout of an acceleration magnitude
\end{definition}
\begin{proof}
  This is the declared model definition of acceleration In Meters Per Second Squared.
\end{proof}
\begin{definition}[Collision Body]
  \label{def:physics:phyx-mini-0815:phyxminiproblems-problemphyxmini0815-collisionbody}
  \lean{PhyXMiniProblems.ProblemPhyXMini0815.CollisionBody}
  Separate bodies present immediately before the impact
\end{definition}
\begin{proof}
  This is the declared model definition of Collision Body.
\end{proof}
\begin{definition}[Figure Stage]
  \label{def:physics:phyx-mini-0815:phyxminiproblems-problemphyxmini0815-figurestage}
  \lean{PhyXMiniProblems.ProblemPhyXMini0815.FigureStage}
  Stages labeled in the supplied ballistic-pendulum image
\end{definition}
\begin{proof}
  This is the declared model definition of Figure Stage.
\end{proof}
\begin{definition}[Swing Stage]
  \label{def:physics:phyx-mini-0815:phyxminiproblems-problemphyxmini0815-swingstage}
  \lean{PhyXMiniProblems.ProblemPhyXMini0815.SwingStage}
  The two states used in the post-impact energy balance
\end{definition}
\begin{proof}
  This is the declared model definition of Swing Stage.
\end{proof}
\begin{definition}[Diagram Axis]
  \label{def:physics:phyx-mini-0815:phyxminiproblems-problemphyxmini0815-diagramaxis}
  \lean{PhyXMiniProblems.ProblemPhyXMini0815.DiagramAxis}
  Coordinate axes explicitly drawn in both panels
\end{definition}
\begin{proof}
  This is the declared model definition of Diagram Axis.
\end{proof}
\begin{definition}[Axis Positive Direction]
  \label{def:physics:phyx-mini-0815:phyxminiproblems-problemphyxmini0815-axispositivedirection}
  \lean{PhyXMiniProblems.ProblemPhyXMini0815.AxisPositiveDirection}
  Positive senses of the displayed coordinate axes
\end{definition}
\begin{proof}
  This is the declared model definition of Axis Positive Direction.
\end{proof}
\begin{definition}[Horizontal Direction]
  \label{def:physics:phyx-mini-0815:phyxminiproblems-problemphyxmini0815-horizontaldirection}
  \lean{PhyXMiniProblems.ProblemPhyXMini0815.HorizontalDirection}
  Horizontal directions used by the two velocity arrows
\end{definition}
\begin{proof}
  This is the declared model definition of Horizontal Direction.
\end{proof}
\begin{definition}[Vertical Direction]
  \label{def:physics:phyx-mini-0815:phyxminiproblems-problemphyxmini0815-verticaldirection}
  \lean{PhyXMiniProblems.ProblemPhyXMini0815.VerticalDirection}
  Vertical direction of the measured pendulum rise
\end{definition}
\begin{proof}
  This is the declared model definition of Vertical Direction.
\end{proof}
\begin{definition}[Collision Outcome]
  \label{def:physics:phyx-mini-0815:phyxminiproblems-problemphyxmini0815-collisionoutcome}
  \lean{PhyXMiniProblems.ProblemPhyXMini0815.CollisionOutcome}
  Qualitative collision outcome
\end{definition}
\begin{proof}
  This is the declared model definition of Collision Outcome.
\end{proof}
\begin{definition}[Pendulum Support]
  \label{def:physics:phyx-mini-0815:phyxminiproblems-problemphyxmini0815-pendulumsupport}
  \lean{PhyXMiniProblems.ProblemPhyXMini0815.PendulumSupport}
  Support geometry stated in the problem and shown in the image
\end{definition}
\begin{proof}
  This is the declared model definition of Pendulum Support.
\end{proof}
\begin{definition}[Collision Impulse Model]
  \label{def:physics:phyx-mini-0815:phyxminiproblems-problemphyxmini0815-collisionimpulsemodel}
  \lean{PhyXMiniProblems.ProblemPhyXMini0815.CollisionImpulseModel}
  Impact-scale external-impulse idealization used for x-momentum balance
\end{definition}
\begin{proof}
  This is the declared model definition of Collision Impulse Model.
\end{proof}
\begin{definition}[Figure Object]
  \label{def:physics:phyx-mini-0815:phyxminiproblems-problemphyxmini0815-figureobject}
  \lean{PhyXMiniProblems.ProblemPhyXMini0815.FigureObject}
  Objects and construction lines visible in image 815.png
\end{definition}
\begin{proof}
  This is the declared model definition of Figure Object.
\end{proof}
\begin{definition}[Figure Label]
  \label{def:physics:phyx-mini-0815:phyxminiproblems-problemphyxmini0815-figurelabel}
  \lean{PhyXMiniProblems.ProblemPhyXMini0815.FigureLabel}
  Symbolic and textual labels printed in image 815.png
\end{definition}
\begin{proof}
  This is the declared model definition of Figure Label.
\end{proof}
\begin{definition}[Ballistic Pendulum Figure]
  \label{def:physics:phyx-mini-0815:phyxminiproblems-problemphyxmini0815-ballisticpendulumfigure}
  \lean{PhyXMiniProblems.ProblemPhyXMini0815.BallisticPendulumFigure}
  \uses{def:physics:phyx-mini-0815:phyxminiproblems-problemphyxmini0815-figurestage, def:physics:phyx-mini-0815:phyxminiproblems-problemphyxmini0815-diagramaxis, def:physics:phyx-mini-0815:phyxminiproblems-problemphyxmini0815-axispositivedirection, def:physics:phyx-mini-0815:phyxminiproblems-problemphyxmini0815-horizontaldirection, def:physics:phyx-mini-0815:phyxminiproblems-problemphyxmini0815-figureobject, def:physics:phyx-mini-0815:phyxminiproblems-problemphyxmini0815-figurelabel}
  Literal qualitative transcription of the primary raster. No field contains the requested initial-speed formula or a numerical value for an unknown
\end{definition}
\begin{proof}
  This is the declared model definition of Ballistic Pendulum Figure.
\end{proof}
\begin{definition}[Ballistic Pendulum Setup]
  \label{def:physics:phyx-mini-0815:phyxminiproblems-problemphyxmini0815-ballisticpendulumsetup}
  \lean{PhyXMiniProblems.ProblemPhyXMini0815.BallisticPendulumSetup}
  \uses{def:physics:phyx-mini-0815:phyxminiproblems-problemphyxmini0815-massquantity, def:physics:phyx-mini-0815:phyxminiproblems-problemphyxmini0815-lengthquantity, def:physics:phyx-mini-0815:phyxminiproblems-problemphyxmini0815-speedquantity, def:physics:phyx-mini-0815:phyxminiproblems-problemphyxmini0815-accelerationquantity, def:physics:phyx-mini-0815:phyxminiproblems-problemphyxmini0815-horizontalmomentumquantity, def:physics:phyx-mini-0815:phyxminiproblems-problemphyxmini0815-energyquantity, def:physics:phyx-mini-0815:phyxminiproblems-problemphyxmini0815-collisionbody, def:physics:phyx-mini-0815:phyxminiproblems-problemphyxmini0815-swingstage, def:physics:phyx-mini-0815:phyxminiproblems-problemphyxmini0815-horizontaldirection, def:physics:phyx-mini-0815:phyxminiproblems-problemphyxmini0815-verticaldirection, def:physics:phyx-mini-0815:phyxminiproblems-problemphyxmini0815-collisionoutcome, def:physics:phyx-mini-0815:phyxminiproblems-problemphyxmini0815-pendulumsupport, def:physics:phyx-mini-0815:phyxminiproblems-problemphyxmini0815-collisionimpulsemodel, def:physics:phyx-mini-0815:phyxminiproblems-problemphyxmini0815-ballisticpendulumfigure}
  Independent physical quantities for the two-stage experiment. In particular, the bullet's initial speed is an unconstrained physical speed; it is not defined from m\_B, m\_W, g, h, or an answer choice
\end{definition}
\begin{proof}
  This is the declared model definition of Ballistic Pendulum Setup.
\end{proof}
\begin{definition}[Matches Stated Scenario]
  \label{def:physics:phyx-mini-0815:phyxminiproblems-problemphyxmini0815-matchesstatedscenario}
  \lean{PhyXMiniProblems.ProblemPhyXMini0815.MatchesStatedScenario}
  \uses{def:physics:phyx-mini-0815:phyxminiproblems-problemphyxmini0815-speedreadout, def:physics:phyx-mini-0815:phyxminiproblems-problemphyxmini0815-ballisticpendulumsetup}
  Qualitative scenario assumptions and the two zero-speed boundary states
\end{definition}
\begin{proof}
  This is the declared model definition of Matches Stated Scenario.
\end{proof}
\begin{definition}[Matches Primary Figure]
  \label{def:physics:phyx-mini-0815:phyxminiproblems-problemphyxmini0815-matchesprimaryfigure}
  \lean{PhyXMiniProblems.ProblemPhyXMini0815.MatchesPrimaryFigure}
  \uses{def:physics:phyx-mini-0815:phyxminiproblems-problemphyxmini0815-figurestage, def:physics:phyx-mini-0815:phyxminiproblems-problemphyxmini0815-diagramaxis, def:physics:phyx-mini-0815:phyxminiproblems-problemphyxmini0815-figurelabel, def:physics:phyx-mini-0815:phyxminiproblems-problemphyxmini0815-ballisticpendulumsetup}
  Facts read directly from the axes, labels, arrows, and geometry in image 815
\end{definition}
\begin{proof}
  This is the declared model definition of Matches Primary Figure.
\end{proof}
\begin{definition}[Has Physical Parameters]
  \label{def:physics:phyx-mini-0815:phyxminiproblems-problemphyxmini0815-hasphysicalparameters}
  \lean{PhyXMiniProblems.ProblemPhyXMini0815.HasPhysicalParameters}
  \uses{def:physics:phyx-mini-0815:phyxminiproblems-problemphyxmini0815-massinkilograms, def:physics:phyx-mini-0815:phyxminiproblems-problemphyxmini0815-lengthinmeters, def:physics:phyx-mini-0815:phyxminiproblems-problemphyxmini0815-speedinmeterspersecond, def:physics:phyx-mini-0815:phyxminiproblems-problemphyxmini0815-accelerationinmeterspersecondsquared, def:physics:phyx-mini-0815:phyxminiproblems-problemphyxmini0815-collisionbody, def:physics:phyx-mini-0815:phyxminiproblems-problemphyxmini0815-ballisticpendulumsetup}
  Positivity and nondegeneracy conditions selecting the physical branches
\end{definition}
\begin{proof}
  This is the declared model definition of Has Physical Parameters.
\end{proof}
\begin{definition}[Satisfies Sticking Collision Momentum Laws]
  \label{def:physics:phyx-mini-0815:phyxminiproblems-problemphyxmini0815-satisfiesstickingcollisionmomentumlaws}
  \lean{PhyXMiniProblems.ProblemPhyXMini0815.SatisfiesStickingCollisionMomentumLaws}
  \uses{def:physics:phyx-mini-0815:phyxminiproblems-problemphyxmini0815-massreadout, def:physics:phyx-mini-0815:phyxminiproblems-problemphyxmini0815-speedreadout, def:physics:phyx-mini-0815:phyxminiproblems-problemphyxmini0815-horizontalmomentumreadout, def:physics:phyx-mini-0815:phyxminiproblems-problemphyxmini0815-ballisticpendulumsetup}
  Mass additivity and one-dimensional x-momentum conservation during the short sticking collision. Momentum is stated in every coherent choice of mass, length, and time units. No field assigns the unknown bullet speed
\end{definition}
\begin{proof}
  This is the declared model definition of Satisfies Sticking Collision Momentum Laws.
\end{proof}
\begin{definition}[Satisfies Conservative Pendulum Swing Laws]
  \label{def:physics:phyx-mini-0815:phyxminiproblems-problemphyxmini0815-satisfiesconservativependulumswinglaws}
  \lean{PhyXMiniProblems.ProblemPhyXMini0815.SatisfiesConservativePendulumSwingLaws}
  \uses{def:physics:phyx-mini-0815:phyxminiproblems-problemphyxmini0815-massreadout, def:physics:phyx-mini-0815:phyxminiproblems-problemphyxmini0815-lengthreadout, def:physics:phyx-mini-0815:phyxminiproblems-problemphyxmini0815-speedreadout, def:physics:phyx-mini-0815:phyxminiproblems-problemphyxmini0815-accelerationreadout, def:physics:phyx-mini-0815:phyxminiproblems-problemphyxmini0815-energyreadout, def:physics:phyx-mini-0815:phyxminiproblems-problemphyxmini0815-swingstage, def:physics:phyx-mini-0815:phyxminiproblems-problemphyxmini0815-ballisticpendulumsetup}
  Kinetic energy, gravitational potential-energy gain, and mechanical-energy conservation for the compound body from immediately after impact to the top of the swing. Energy conservation is not asserted across the inelastic collision
\end{definition}
\begin{proof}
  This is the declared model definition of Satisfies Conservative Pendulum Swing Laws.
\end{proof}
\begin{lemma}[post Impact Speed eq sqrt two g h]
  \label{lem:physics:phyx-mini-0815:phyxminiproblems-problemphyxmini0815-postimpactspeed-eq-sqrt-two-g-h}
  \lean{PhyXMiniProblems.ProblemPhyXMini0815.postImpactSpeed_eq_sqrt_two_g_h}
  \uses{def:physics:phyx-mini-0815:phyxminiproblems-problemphyxmini0815-lengthinmeters, def:physics:phyx-mini-0815:phyxminiproblems-problemphyxmini0815-speedinmeterspersecond, def:physics:phyx-mini-0815:phyxminiproblems-problemphyxmini0815-accelerationinmeterspersecondsquared, def:physics:phyx-mini-0815:phyxminiproblems-problemphyxmini0815-ballisticpendulumsetup, def:physics:phyx-mini-0815:phyxminiproblems-problemphyxmini0815-matchesstatedscenario, def:physics:phyx-mini-0815:phyxminiproblems-problemphyxmini0815-hasphysicalparameters, def:physics:phyx-mini-0815:phyxminiproblems-problemphyxmini0815-satisfiesconservativependulumswinglaws}
  Energy conservation during the swing selects the nonnegative square-root branch v\_2 = sqrt (2 g h). This is an intermediate conclusion, not an assumption of the collision model
\end{lemma}
\begin{proof}
  The conclusion follows from the governing relations and figure readouts named in the statement of post Impact Speed eq sqrt two g h.
\end{proof}
\begin{lemma}[bullet Initial Speed eq mass Ratio mul post Impact Speed]
  \label{lem:physics:phyx-mini-0815:phyxminiproblems-problemphyxmini0815-bulletinitialspeed-eq-massratio-mul-postimpactspeed}
  \lean{PhyXMiniProblems.ProblemPhyXMini0815.bulletInitialSpeed_eq_massRatio_mul_postImpactSpeed}
  \uses{def:physics:phyx-mini-0815:phyxminiproblems-problemphyxmini0815-massinkilograms, def:physics:phyx-mini-0815:phyxminiproblems-problemphyxmini0815-speedinmeterspersecond, def:physics:phyx-mini-0815:phyxminiproblems-problemphyxmini0815-ballisticpendulumsetup, def:physics:phyx-mini-0815:phyxminiproblems-problemphyxmini0815-matchesstatedscenario, def:physics:phyx-mini-0815:phyxminiproblems-problemphyxmini0815-hasphysicalparameters, def:physics:phyx-mini-0815:phyxminiproblems-problemphyxmini0815-satisfiesstickingcollisionmomentumlaws}
  Horizontal momentum conservation across the embedding collision gives the mass-ratio relation between v\_1 and v\_2
\end{lemma}
\begin{proof}
  The conclusion follows from the governing relations and figure readouts named in the statement of bullet Initial Speed eq mass Ratio mul post Impact Speed.
\end{proof}
\begin{definition}[Answer Choice]
  \label{def:physics:phyx-mini-0815:phyxminiproblems-problemphyxmini0815-answerchoice}
  \lean{PhyXMiniProblems.ProblemPhyXMini0815.AnswerChoice}
  Labels attached to the four formulas printed with the problem
\end{definition}
\begin{proof}
  This is the declared model definition of Answer Choice.
\end{proof}
\begin{definition}[displayed Initial Speed Formula In Meters Per Second]
  \label{def:physics:phyx-mini-0815:phyxminiproblems-problemphyxmini0815-displayedinitialspeedformulainmeterspersecond}
  \lean{PhyXMiniProblems.ProblemPhyXMini0815.displayedInitialSpeedFormulaInMetersPerSecond}
  \uses{def:physics:phyx-mini-0815:phyxminiproblems-problemphyxmini0815-massinkilograms, def:physics:phyx-mini-0815:phyxminiproblems-problemphyxmini0815-lengthinmeters, def:physics:phyx-mini-0815:phyxminiproblems-problemphyxmini0815-accelerationinmeterspersecondsquared, def:physics:phyx-mini-0815:phyxminiproblems-problemphyxmini0815-ballisticpendulumsetup, def:physics:phyx-mini-0815:phyxminiproblems-problemphyxmini0815-answerchoice}
  The speed expression printed beside each answer label, in metres per second
\end{definition}
\begin{proof}
  This is the declared model definition of displayed Initial Speed Formula In Meters Per Second.
\end{proof}
\begin{definition}[recorded Dataset Answer]
  \label{def:physics:phyx-mini-0815:phyxminiproblems-problemphyxmini0815-recordeddatasetanswer}
  \lean{PhyXMiniProblems.ProblemPhyXMini0815.recordedDatasetAnswer}
  \uses{def:physics:phyx-mini-0815:phyxminiproblems-problemphyxmini0815-answerchoice}
  The answer label recorded in the supplied dataset metadata
\end{definition}
\begin{proof}
  This is the declared model definition of recorded Dataset Answer.
\end{proof}
% --- Archon declaration topology end ---
% --- Archon physics formalization source end ---
